\section{Algebra delle Matrici}

\subsection{Operazioni con Matrici}

\begin{exercisebox}[title={Esercizio 1: Prodotto Righe per Colonne}]
Date le matrici:
$$ A = \begin{pmatrix} 1 & 2 & 0 \\ -1 & 3 & 1 \end{pmatrix}, \quad B = \begin{pmatrix} 2 & -1 \\ 0 & 1 \\ 3 & 4 \end{pmatrix} $$
Calcolare, se possibile, i prodotti $AB$ e $BA$.
\end{exercisebox}

\begin{solbox}
\textbf{Calcolo di $AB$:}
Scriviamo le matrici da moltiplicare:
$$ AB = \begin{pmatrix} 1 & 2 & 0 \\ -1 & 3 & 1 \end{pmatrix} \begin{pmatrix} 2 & -1 \\ 0 & 1 \\ 3 & 4 \end{pmatrix} $$
La matrice prodotto $AB$ ha dimensione $2 \times 2$. Calcoliamo i termini $c_{ij}$:
\begin{itemize}
    \item $c_{11} = 1(2) + 2(0) + 0(3) = 2$
    \item $c_{12} = 1(-1) + 2(1) + 0(4) = 1$
    \item $c_{21} = -1(2) + 3(0) + 1(3) = 1$
    \item $c_{22} = -1(-1) + 3(1) + 1(4) = 8$
\end{itemize}
Risultato:
$$ AB = \begin{pmatrix} 2 & 1 \\ 1 & 8 \end{pmatrix} $$

\textbf{Calcolo di $BA$:}
Scriviamo le matrici in ordine inverso:
$$ BA = \begin{pmatrix} 2 & -1 \\ 0 & 1 \\ 3 & 4 \end{pmatrix} \begin{pmatrix} 1 & 2 & 0 \\ -1 & 3 & 1 \end{pmatrix} $$
La matrice prodotto $BA$ ha dimensione $3 \times 3$. Calcoliamo i termini:
\begin{itemize}
    \item Riga 1: $2(1)+(-1)(-1)=3, \quad 2(2)+(-1)(3)=1, \quad 2(0)+(-1)(1)=-1$
    \item Riga 2: $0(1)+1(-1)=-1, \quad 0(2)+1(3)=3, \quad 0(0)+1(1)=1$
    \item Riga 3: $3(1)+4(-1)=-1, \quad 3(2)+4(3)=18, \quad 3(0)+4(1)=4$
\end{itemize}
Risultato:
$$ BA = \begin{pmatrix} 3 & 1 & -1 \\ -1 & 3 & 1 \\ -1 & 18 & 4 \end{pmatrix} $$
\end{solbox}

\begin{exercisebox}[title={Esercizio 2: Inversa con Gauss-Jordan}]
Calcolare l'inversa della matrice:
$$ A = \begin{pmatrix} 1 & 0 & 1 \\ 1 & 1 & 2 \\ 0 & 1 & 2 \end{pmatrix} $$
\end{exercisebox}

\begin{solbox}
Scriviamo la matrice completa $(A|I)$:
$$ \left(\begin{array}{ccc|ccc} 
1 & 0 & 1 & 1 & 0 & 0 \\ 
1 & 1 & 2 & 0 & 1 & 0 \\ 
0 & 1 & 2 & 0 & 0 & 1 
\end{array}\right) $$
Applichiamo l'eliminazione di Gauss-Jordan.
1. $R_2 \to R_2 - R_1$:
$$ \left(\begin{array}{ccc|ccc} 1 & 0 & 1 & 1 & 0 & 0 \\ 0 & 1 & 1 & -1 & 1 & 0 \\ 0 & 1 & 2 & 0 & 0 & 1 \end{array}\right) $$
2. $R_3 \to R_3 - R_2$:
$$ \left(\begin{array}{ccc|ccc} 1 & 0 & 1 & 1 & 0 & 0 \\ 0 & 1 & 1 & -1 & 1 & 0 \\ 0 & 0 & 1 & 1 & -1 & 1 \end{array}\right) $$
3. $R_2 \to R_2 - R_3$ (Risalita):
$$ \left(\begin{array}{ccc|ccc} 1 & 0 & 1 & 1 & 0 & 0 \\ 0 & 1 & 0 & -2 & 2 & -1 \\ 0 & 0 & 1 & 1 & -1 & 1 \end{array}\right) $$
4. $R_1 \to R_1 - R_3$ (Risalita):
$$ \left(\begin{array}{ccc|ccc} 1 & 0 & 0 & 0 & 1 & -1 \\ 0 & 1 & 0 & -2 & 2 & -1 \\ 0 & 0 & 1 & 1 & -1 & 1 \end{array}\right) $$
La matrice a sinistra è l'identità. L'inversa è la matrice a destra:
$$ A^{-1} = \begin{pmatrix} 0 & 1 & -1 \\ -2 & 2 & -1 \\ 1 & -1 & 1 \end{pmatrix} $$
\end{solbox}

\begin{exercisebox}[title={Esercizio 3: Potenze di Matrici}]
Data $J = \begin{pmatrix} 0 & 1 \\ 0 & 0 \end{pmatrix}$, calcolare $J^2$ e $J^3$.
\end{exercisebox}

\begin{solbox}
Calcoliamo il quadrato $J^2$:
$$ J^2 = \begin{pmatrix} 0 & 1 \\ 0 & 0 \end{pmatrix} \begin{pmatrix} 0 & 1 \\ 0 & 0 \end{pmatrix} = \begin{pmatrix} 0 & 0 \\ 0 & 0 \end{pmatrix} = \mathbf{0} $$
Dato che $J^2$ è la matrice nulla, segue immediatamente che:
$$ J^3 = J^2 \cdot J = \mathbf{0} \cdot J = \mathbf{0} $$
La matrice $J$ è nilpotente di indice 2.
\end{solbox}
